% !PRINTED 620
\weberpage{620}

\begin{center}
\S. 170.

\gesperrt{Minimal basis of the field $\Omega_m$.}
\end{center}

The theorems that were proved in the preceding section lead us to the following theorem:

II. \gesperrt{The numbers}
\weq{1}{1, r, r^2, \ldots, r^{\mu-1}}
\gesperrt{form a basis of the system $\mathfrak{o}$ of integers in $\Omega_m$.}

To show this it suffices, by \S.\,145, to prove that the integer in $\Omega_m$
\weq{2}{\omega = x_0 + x_1 r + x_2 r^2 + \cdots + x_{\mu-1} r^{\mu-1},}
in which $x_0, x_1, \ldots, x_{\mu-1}$ are integral rational numbers, can be divisible by a rational prime $p$ only when the numbers $x_0, x_1, \ldots, x_{\mu-1}$ are all divisible by $p$.

Let us assume, then, that $\omega$ is divisible by $p$; then all the numbers $\omega_n$ that arise from $\omega$ by one of the $\mu$ substitutions $(r, r^n)$ are also divisible by $p$, and thus from (2) we obtain a system of $\mu$ equations which are linear with respect to the $x_i$:
\weq{3}{\omega_n = x_0 + x_1 r^n + x_2 r^{2n} + \cdots + x_{\mu-1} r^{(\mu-1)n}.}

But the determinant of this system, that is, the determinant formed from the $\mu$ rows
\[ 1, r^n, r^{2n}, \ldots, r^{(\mu-1)n} \]
is, by Bd. I, \S.\,46, equal to the square root $\sqrt{\varDelta}$, and if we now solve the system of equations (3), it follows that all the rational numbers $\varDelta x_i$ must be divisible by $p$.

If now $p$ is not $=q$, then $\varDelta$ is not divisible by $p$, and all the $x_i$ must be divisible by $p$.

But if $p = q$, we set
\[ f(t) = x_0 + x_1 t + x_2 t^2 + \cdots + x_{\mu-1} t^{\mu-1}, \]
so that $\omega_n = f(r^n)$, and in this we set $r = 1 - \sigma$ [\S.\,169, (6)]. Then, according to the \gesperrt{Taylor} development:
\weq{4}{\begin{aligned}
\omega &= f(1-\sigma) = f(1) - \sigma f'(1) + \dfrac{\sigma^2}{1.2} f''(1) - \cdots\\
&\qquad - \dfrac{\sigma^{\mu-1}}{\Pi(\mu-1)} f^{(\mu-1)}(1);
\end{aligned}}
%CONT>>
